% ============================================================================
% RELATED: DETERMINISTIC ENCODING AND ARITHMETIC CODING
% ============================================================================
\subsection{Deterministic Encoding and Collision Properties}
\label{sec:related_godel_arithmetic}

The system encodes input bytes into Value frames through a deterministic
mapping that assigns each byte sequence a fixed binary representation.
This subsection places this encoding in the context of two classical
results: G\"{o}del numbering and arithmetic coding.

\paragraph{Relationship to G\"{o}del Numbering.}
G\"{o}del's 1931 encoding assigns every formula of a formal language a
unique natural number via products of prime powers.  The Fundamental
Theorem of Arithmetic guarantees that this mapping is injective: distinct
formulas receive distinct numbers, and the original formula can be
recovered from its number.  The Value encoding shares the structural
goal of providing a deterministic, reproducible mapping from symbolic
input to a fixed binary representation, though it operates at the byte
level rather than over formal language symbols.

The key property inherited from the G\"{o}del construction is
\emph{determinism}: the same input byte sequence always produces the
same Value content, and distinct inputs that differ in at least one byte
position produce Values that differ in the corresponding data-region
words.  This guarantee holds without learning or randomized hashing; it
is a consequence of the encoding rule itself.

\paragraph{Collision and Deduplication.}
When two distinct input sequences produce Values that share a common
prefix in their data regions, the shared structure can be exploited for
storage efficiency.  This is analogous to the principle underlying
arithmetic coding \citep{witten1987arithmetic}, in which symbols are
mapped to sub-intervals of $[0, 1)$ proportional to their empirical
probability, achieving entropy-rate compression.

In the Value architecture, deduplication arises naturally when multiple
inputs map to identical data-region contents.  Because the mapping is
deterministic, identical byte sequences always produce identical Values,
and a storage layer that indexes Values by their content can coalesce
duplicates without additional computation.  The degree of deduplication
depends on the redundancy structure of the input data: highly repetitive
corpora (such as natural language, where Zipf's law governs token
frequencies) produce more collisions than uniformly random input.

\paragraph{Implications for Storage Scaling.}
This collision property means that the effective memory footprint of a
collection of Values scales with the information content of the data
rather than its raw volume.  The compression ratio is bounded below by
the Shannon entropy of the input distribution---the same theoretical
floor that governs arithmetic coding.  Empirical measurements of
deduplication rates across different data domains are reported in the
scaling experiments.
